% Esquisse d'un programme (cote 119, copy 1): archive cote 119, site page 38 (PDF page 39), typescript page 36.
% typescript. Draft compared: EsquisseFr.pdf (text layer):1300-1352.
% AI-transcribed by Claude Opus 5.5 (claude-opus-5-5, adaptive thinking, effort high), Message Batches API, batch esq1 (msgbatch_01LNsogFyF3y6KSp9UbtJtBk), 2026-11-12; not yet checked by eye.
\scanpage{119/38}
\mspage{36}

intervenant dans la situation, la théorie de dévissage à laquelle
j'étais parvenue dans un contexte tout différent - avec la tâche
supplémentaire, par la suite, de dégager un théorème de comparai-%
son général, sur le modèle des théorèmes bien connus, pour com-%
parer les invariants obtenus (notamment les types d'homotopie de
voisinages tubulaires divers) dans le cadre transcendant, et dans
le cadre schématique. J'ai pu me convaincre qu'un tel formalisme
de dévissage avait bel et bien un sens dans le contexte (dit
``abstrait'' !) des topos généraux, ou tout au moins des topos noe-%
thériens (comme ceux qui s'introduisent ici), via une notion conve-%
nable de \uline{voisinage tubulaire canonique d'un sous-topos} dans un
topos ambiant. Une fois cette notion acquise, avec certaines pro-%
priétés formelles simples, la description du ``dévissage'' d'un
topos stratifié est considérablement plus simple même dans ce
cadre, que dans le cadre topologique (modéré). Il est vrai que
là aussi il y a un travail de fondements à faire, notamment pour
la notion même de voisinage tubulaire d'un sous-topos - et il est
étonnant d'ailleurs que ce travail (pour autant que je sache)
n'ait toujours pas été fait, c'est-à-dire que personne (depuis plus
de vingt ans qu'il existe un contexte de topologies étale) ne
semble en avoir eu besoin ; un signe sûrement que la compréhension
de la structure topologique des schémas n'a pas tellement pro-%
gressées depuis le travail d'Artin-Mazur ...

Une fois accompli le double travail de dégrossissage
(plus ou moins heuristique) autour de la notion de dévissage d'un
espace ou d'un topos stratifié, qui a été une étape cruciale dans
ma compréhension des multiplicités modulaires, il est d'ailleurs
apparu que pour les besoins de ces dernières, on peut sans doute
court-circuiter au moins une bonne partie de cette théorie par des
arguments géométriques directs. Il n'en reste pas moins que pour
moi, le formalisme de dévissage auquel je suis parvenu a fait ses
preuves d'utilité et de cohérence, indépendamment de toute ques-%
tion sur les fondements les plus adéquats qui permettent de lui
donner tout son sens.

\margin{38}
